% ================================================================
%  SECTION 4 -- rev 2, after reading ch. 5 of RL_for_PDEs__Ali_Tayyar.pdf
%  Target: 2.00 pages (rebalanced from 1.50; see note in chat).
%  Sources: thesis ch. 5 (pp. 27-38) + current notebook code.
%  Where the two disagree, the CODE is authoritative and the disagreement
%  is flagged inline.
% ================================================================

\section{Placement as a Markov Decision Process}
\label{sec:mdp}

Once the neighbour set is chosen, nothing else is free: the weights follow
from the moment conditions of \cref{sec:recon}, and the error is bounded by
$\normone{\bw}\, C h^2$. The only open choice is which points to evaluate,
and in every classical method that choice is made in advance --- by a grid,
or by a refinement rule written by hand. We give it to an agent instead.

An action is then a decision to spend one evaluation. A single episode
answers a single query $\zstar = (x^*, t^*)$: it spends evaluations until the
operator admits an admissible stencil at $\zstar$, and is scored on the error
it achieves there and on how many evaluations that took.

The agent works through $N$ \emph{walkers}, initialised on the initial
condition, each carrying a solution value. The environment maintains a
visited set $\Vset$ holding every evaluated point --- the initial condition,
the boundary values, and every accepted solve. At each step the agent picks
one walker and one displacement. The proposal is checked against the domain
bounds and against $\Vset$ for collision; if it survives, a solve is
attempted there using neighbours from $\Vset$. On success the point joins
$\Vset$ and the walker moves; on failure the walker stays and the step is
recorded as a rejection.

\subsection{State}
\label{sec:mdp_state}

The observation is the walker configuration expressed relative to the query,
\begin{equation}
  s = \Bigl(\tfrac{x_1 - x^*}{\Delta x}, \ldots, \tfrac{x_N - x^*}{\Delta x},\;
            \tfrac{t_1 - t^*}{\Delta t}, \ldots, \tfrac{t_N - t^*}{\Delta t},\;
            \tfrac{t^*}{\Delta t},\; \tfrac{x^*}{\Delta x} \Bigr)
  \in \R^{2N+2}.
  \label{eq:state}
\end{equation}
Every entry is a displacement in grid cells. This matters because the
quantities that govern admissibility --- two-sided coverage, collision,
conditioning --- all live at cell resolution and do not rescale with the
horizon. An earlier absolute-coordinate state carried no information about
the discretisation scale: nothing in the pair $(0.05, 0.1)$ tells the policy
whether the query is five time steps away or five hundred. A fractional
parameterisation $t/t^*$ was considered and rejected for the mirror-image
reason: it makes ``one cell from the target'' mean $0.02$ at one horizon and
$0.1$ at another, blinding the agent precisely where the problem is not
scale-free.

Two omissions are deliberate. The solution values $u_i$ are not observed, so
the policy chooses geometry without seeing the field it reconstructs. More
consequentially, $\Vset$ is not observed either, even though it is $\Vset$
--- not the walker positions --- that determines whether a stencil is
admissible and how accurate the prediction is. The problem is therefore
partially observed, and the walker configuration is a summary statistic
rather than a state. We keep it because the walkers generate $\Vset$, so
their configuration is informative about it, and because a fixed-size
observation is what allows one policy to serve queries at different
locations; \cref{sec:conclusion} returns to the cost of this choice.

\subsection{Actions and admissibility masking}
\label{sec:mdp_action}

A walker steps to a neighbouring cell,
\begin{equation}
  \Aset = \{-1, 0, +1\} \times \{+1, +2\}
  \quad\text{in units of } (\Delta x, \Delta t),
  \label{eq:actions}
\end{equation}
so evaluations always advance in time. Displacements with $\delta t = 0$ were
removed after they were observed to consume roughly half of all actions:
walkers advance monotonically in time and $\Vset$ already contains their own
trail, so a move at fixed $t$ almost always lands on an occupied cell. The
cost of removing them is recorded --- no walker can reposition laterally
without also advancing --- so lateral coverage of a time level must come
from having several walkers at different $x$ rather than from any one walker
spreading out.

A joint action selects both a walker and its move, giving a discrete space of
size $N \cdot |\Aset|$, decoded as $i = \lfloor a / |\Aset| \rfloor$,
$j = a \bmod |\Aset|$. The flattened form matters. Factoring the policy into
independent walker and displacement heads,
$\pi(i,j \mid s) = \pi^{\mathrm{w}}(i \mid s)\,\pi^{\mathrm{m}}(j \mid s)$,
is a rank-one outer product, and what lives in the missing dimensions is
exactly the correlation between which walker moves and where it goes. An
optimal policy of the form ``if walker 3, go left; if walker 7, go right''
is not expressible: putting mass on walkers $\{3,7\}$ and moves
$\{\text{left},\text{right}\}$ necessarily also puts mass on
$(3, \text{right})$ and $(7, \text{left})$. Under deterministic rollout the
arg-max is assembled from two near-uniform marginals and lands on a pair the
policy never intended, which is what opened a gap between sampled and greedy
performance. A single softmax over the joint table is full-rank and closes
the gap by construction.

Not every action is available. A proposal is masked if it leaves the domain
or the ceiling $t^* + R\,\Delta t$; if its target cell is already in $\Vset$;
or if a previous attempt there failed to admit a stencil, in which case the
cell is blocked until $\Vset$ has grown by a further $k$ points, on the
reasoning that the available stencil will have changed by then. The first
two are computable before acting; the third is not, since whether a solve
succeeds is only discovered by attempting it. Masking is therefore not a
static function of position --- it accumulates the episode's own history of
failed solves.

Masking was the single most consequential implementation change. A rejected
proposal leaves \cref{eq:state} unchanged, since the walker does not move. A
policy is a function of the observation, so having selected an action that is
rejected it is presented with the identical observation and selects the
identical action indefinitely; two-action cycles were also observed.
Stochastic sampling escapes such cycles only by wasting a large fraction of
the budget. Maskable PPO \CITE{huang2022masking} renormalises over admissible
actions and makes the pathology unreachable rather than merely improbable.

\subsection{Termination}
\label{sec:mdp_term}

A box of half-width $R$ cells is placed around the query,
$B_R(\zstar) = \{(x,t) : |x - x^*| \le R\,\Delta x,\; |t - t^*| \le R\,\Delta t\}$,
and the episode terminates when $\zstar$ admits an admissible stencil drawn
from $B_R(\zstar)$, or when a budget of $\Kmax$ steps is exhausted. No walker
need land on $\zstar$. Walkers may rise above the query, up to
$t^* + R\,\Delta t$, so the box can be populated from both sides in time; the
final solve is correspondingly non-causal, whereas every intermediate solve
during the rollout uses causal neighbours only.

The stopping rule deliberately asks the same question as the scoring rule
--- \emph{does an admissible stencil exist at $\zstar$} --- rather than the
cheaper proxy \emph{are $n$ evaluated points nearby}. The distinction is not
cosmetic. Under the proxy an episode can satisfy termination and then fail
its final solve, receiving the worst accuracy penalty regardless of anything
the agent did; the reward is then constant across those episodes and carries
no gradient. \TODO{quote the fraction of episodes affected under the proxy
rule --- roughly 0.8 in the diagnosis. Verify before quoting.}

\subsection{Reward}
\label{sec:mdp_reward}

The return is terminal, and both branches are scored by the same expression.
On termination the operator is applied at $\zstar$ and the error
$e = |\ux - \uhat|$ is compressed onto a bounded scale,
\begin{equation}
  \enorm = \mathrm{clip}\!\left(
    \frac{\log_{10}(e / \etol)}{3},\; 0,\; 1 \right),
  \qquad
  r = -\,w_{\mathrm{err}}\,\enorm \;-\; w_{\mathrm{time}}\,\frac{K}{\Kmax},
  \label{eq:reward}
\end{equation}
with $\enorm = 1$ when the final solve admits no stencil. The logarithmic
form replaced a linear saturation $\min(e/\etol, 1)$ that resolved only one
decade; it grades three. It was itself preceded by an unnormalised reward
$-e - \lambda K$, which failed in an instructive way: the requirement that
reaching always beat timing out could not be met at every horizon by any
choice of constants, because accumulated error eventually made the reached
branch worse than the timeout branch, and the only repair available by tuning
was to shrink the accuracy weight --- that is, to declare poor accuracy
acceptable. Normalisation was adopted because the ordering had to hold
structurally rather than by calibration.

We use $\gamma = 1$. A discount would assert that evaluations placed early
are worth more than those placed late, which is not a preference we hold: an
evaluation costs the same wherever it occurs, and the accuracy at $\zstar$ is
the same quantity regardless of when the contributing points were placed.

A potential-based shaping term densifies an otherwise terminal signal. With
$\phi(s) = -c_{\mathrm{shape}}\, N^{-1} \sum_i d_i / d_0$, where $d_i$ is the
normalised walker--query distance and $d_0$ its value at reset, and
$\phi \equiv 0$ at terminal states, the contributions telescope to
$-\phi(s^{(0)})$ --- a constant depending only on the initial state. The
shaping therefore cannot change the optimal policy \CITE{ng1999policy} and
acts purely on the variance of the return. Two consequences are worth
stating. An earlier version used $\min_i d_i$ rather than the mean, under
which a move by any walker other than the current closest leaves $\phi$
unchanged, so $N-1$ of the $N$ walker selections receive no shaping at all;
it also rewards a single walker sprinting ahead, which by
\cref{sec:recon_bound} is precisely the geometry that inflates
$\normone{\bw}$. And since the episode total is a constant proportional to
$c_{\mathrm{shape}}$, returns are not comparable across runs with different
shaping coefficients.

\subsection{Admissibility under exploration}
\label{sec:mdp_admiss}

The reconstruction environment enforces four admissibility conditions:
two-sided spatial coverage, a bound on $\mathrm{cond}(\bA)$, a bound on the
largest weight magnitude, and a lower bound on negative weights. Carried
unchanged into the RL environment, all four together left the walkers barely
able to move: nearly every proposal was rejected, so $\Kmax$ had to be raised
substantially for a walker to reach the query at all, and the reach rate
collapsed. Under a sparse terminal reward this is fatal in a specific way ---
an agent that rarely reaches the query never discovers that a reward exists
there, and further training does not help. The two weight constraints were
therefore disabled and the conditioning tests retained.

The cost of that relaxation is exactly quantified by \cref{sec:recon_bound}:
the two disabled constraints are the only bounds the environment had on
$\normone{\bw}$, which is simultaneously the factor in the error bound and
the per-level amplification rate. \Cref{sec:breaks} shows that this is what
sets the long-horizon ceiling.

\subsection{Training}
\label{sec:mdp_policy}

We train with maskable PPO \CITE{schulman2017ppo, sb3contrib} on a
multilayer-perceptron policy over $12$ subprocess environments.
\Cref{tab:hyper} lists the configuration. GPU execution was tested and found
slower than CPU: the policy network is small and the cost is dominated by the
environment step --- a neighbour search and a $3 \times 5$ least-squares
solve --- so transfer overhead outweighs any gain.

\begin{table}[t]
  \centering
  \caption{Environment and training configuration.
  \TODO{must match the run featured in \cref{sec:exp}; values below are the
  current notebook defaults. NOTE: the thesis chapter lists entropy
  coefficient $0.01$, the code uses $0.005$ --- resolve.}}
  \label{tab:hyper}
  \begin{tabular}{@{}llll@{}}
    \toprule
    Environment & & Training & \\
    \midrule
    Walkers $N$              & \NUM{70}      & Algorithm      & MaskablePPO \\
    Bubble radius $R$        & \NUM{2}       & Parallel envs  & \NUM{12} \\
    Stencil size $n$         & \NUM{5}       & Rollout length & \NUM{1024} \\
    Step budget $\Kmax$      & \NUM{5000}    & Batch size     & \NUM{1024} \\
    Tolerance $\etol$        & \NUM{10^{-4}} & Learning rate  & \NUM{2 \times 10^{-4}} \\
    $w_{\mathrm{err}}$       & \NUM{1.0}     & Clip range     & \NUM{0.1} \\
    $w_{\mathrm{time}}$      & \NUM{0.1}     & Entropy coef.  & \NUM{0.005} \\
    $c_{\mathrm{shape}}$     & \NUM{0.3}     & Discount $\gamma$ & \NUM{1.0} \\
    \bottomrule
  \end{tabular}
\end{table}

The environment accepts a schedule of queries, drawn per episode, so that
each query supplies a distinct set of training examples under the relative
state of \cref{eq:state}. The implemented scheduler samples jointly and
uniformly rather than running a sequential curriculum; the transfer argument
that motivated scheduling is really an argument for the sequential form,
which has not been tested. \TODO{state plainly whether the featured run uses
a single fixed query. If it does, $t^*/\Delta t$ and $x^*/\Delta x$ are
constant inputs and the generalisation machinery is inactive --- say so here
rather than let a reviewer find it.}
